% SOURCE_AUTHORITY: sga5_fr_workpass.tex, lines 12072--12115 (LF-normalized unit).
% SOURCE_SHA256: 791F4EFFC5E02832D5D77ED03518C8156D6F07E4C8238B03545DB93D883FBB28
\emph{Demostración del lema.} Sea $u:F\to F'$ un epimorfismo. Hay que demostrar que
\[
\Hom_A(E,F)\longrightarrow\Hom_A(E,F')
\]
es un epimorfismo. Basta demostrar que, para todo grupo divisible $G$,
\[
\Hom_{\mathbb Z}(\Hom_A(E,F'),G)\longrightarrow
\Hom_{\mathbb Z}(\Hom_A(E,F),G)
\]
es un monomorfismo. Ahora bien, el diagrama
\[
\begin{tikzcd}[column sep=large,row sep=large]
\Hom_{\mathbb Z}(\Hom_A(E,F'),G)
 \arrow[r]\arrow[d,"\sim"']&
\Hom_{\mathbb Z}(\Hom_A(E,F),G)\arrow[d,"\sim"]\\
\Hom_{\mathbb Z}(F',G)\otimes_A E
 \arrow[r]&
\Hom_{\mathbb Z}(F,G)\otimes_A E
\end{tikzcd}
\]
es conmutativo, y las flechas verticales son isomorfismos gracias a la hipótesis hecha sobre $E$.

\emph{Demostración de la proposición 5.2.} La implicación a) $\Rightarrow$ b) es inmediata. Para demostrar b) $\Rightarrow$ a), sea
\[
\cdots\longrightarrow P_n\xrightarrow{d_n}\cdots\longrightarrow P_i\longrightarrow 0\longrightarrow\cdots
\]
un complejo isomorfo a $L^\bullet$ en $D(\cC)$, siendo los $P_j$ $A$-módulos proyectivos de tipo finito. Como el complejo $L^\bullet$ tiene dimensión Tor finita y su cohomología es acotada, resulta que existe $n_o\in\mathbb Z$ tal que:

I) La sucesión
\[
\cdots\longrightarrow P_{n_o-p}\xrightarrow{d_{n_o-p}}\cdots\longrightarrow
P_{n_o}\longrightarrow \Ker d_{n_o+1}\longrightarrow0
\]
es exacta.

II) $\Ker d_{n_o+1}$ es plano.

III) El complejo
\[
0\longrightarrow \Ker d_{n_o+1}\longrightarrow P_{n_o+1}\longrightarrow\cdots\longrightarrow P_i\longrightarrow0\longrightarrow\cdots
\]
es isomorfo a $L^\bullet$ en $D(\cC)$.

Ahora bien, $\Ker d_{n_o+1}$ es de presentación finita y, por tanto, proyectivo de tipo finito gracias al lema, lo que demuestra 8.1.
